\appendix

\section{Derivation Details}
\label{app:derivations}

\subsection{Signal Chain: From DAC to Hall Voltage}

Starting from the DAC output voltage $\Vda$:

\begin{enumerate}
\item \textbf{Current:}
$I = \kA \cdot \Vda$, \quad $[\si{\ampere}] = [\si{\ampere/\volt}] \times [\si{\volt}]$.

\item \textbf{Flux (Hopkinson's law):}
$\Phi = \frac{\Nc \cdot I}{\Ra} = \kpole \cdot I$, \quad
$[\si{\weber}] = [\si{\weber/\ampere}] \times [\si{\ampere}]$.

\item \textbf{Magnetic charge:}
$q = \frac{\Phi}{\mu_0}$, \quad
$[\si{\ampere\cdot\metre}] = [\si{\weber}] / [\si{\tesla\cdot\metre/\ampere}]$.

\item \textbf{Coulomb's law:}
$B = \frac{\mu_0}{4\pi} \cdot \frac{q}{r^2}
   = \frac{\mu_0}{4\pi} \cdot \frac{\kpole I}{\mu_0 \, r^2}
   = \frac{\kpole I}{4\pi \, r^2}$.

Note: $\mu_0$ cancels between steps 3 and 4.

\item \textbf{Hall voltage:}
$\Vm = \SH \cdot B = \frac{\SH \cdot \kpole \cdot I}{4\pi \, r^2}$.
\end{enumerate}

Substituting $I = \kA \Vda$ and $r = x + b$ (with $x$, $b$ in micrometers):
\begin{equation}
H(x) = \frac{\Vm}{\Vda} = \frac{\SH \cdot \kpole \cdot \kA}{4\pi} \cdot
\frac{10^{12}}{(x + b)^2} = \frac{a^2}{(x+b)^2},
\tag{A.1}
\end{equation}
where $a^2 = \SH \cdot \kpole \cdot \kA \times 10^{12} / (4\pi)$.

\subsection{Force on a Paramagnetic Bead}

For a paramagnetic bead (radius $r_b$, volume susceptibility $\chi$) in a
non-uniform field $B(x)$:

\begin{enumerate}
\item \textbf{Effective susceptibility} (sphere demagnetization):
$\chitot = \frac{3\chi}{3 + \chi}$.

\item \textbf{Bead volume:}
$\Vbead = \frac{4}{3}\pi r_b^3$.

\item \textbf{Magnetization energy:}
$U = -\frac{\chitot \Vbead}{2\mu_0} B^2$.

\item \textbf{Force:}
$\Fmag = -\frac{\mathrm{d}U}{\mathrm{d}x}
       = \frac{\chitot \Vbead}{\mu_0} B \left|\frac{\mathrm{d}B}{\mathrm{d}x}\right|$.
\end{enumerate}

Substituting the monopole field:
\begin{align}
B(x) &= \frac{\kpole I}{4\pi (x+b)^2}, \tag{A.2} \\
\frac{\mathrm{d}B}{\mathrm{d}x} &= -\frac{2\kpole I}{4\pi (x+b)^3}, \tag{A.3} \\
\Fmag(x) &= \frac{\chitot \Vbead}{\mu_0} \cdot
\frac{2(\kpole I)^2}{(4\pi)^2 (x+b)^5}. \tag{A.4}
\end{align}

\subsection{Thermal Force}

The root-mean-square thermal force on a bead in a viscous medium over a
measurement timescale $\Delta t$ is
\begin{equation}
\Fthermal = \sqrt{\frac{2 k_B T \gamma}{\Delta t}},
\tag{A.5}
\end{equation}
where $\gamma = 6\pi \eta r_b$ is the Stokes drag coefficient,
$\eta = \SI{1e-3}{\pascal\cdot\second}$ for water, and
$k_B = 1.381 \times 10^{-23}$ J/K.

\section{Complete Data Table}
\label{app:data-table}

\Cref{tab:full-data} presents the complete $22 \times 15$ transfer function
magnitude matrix $|H(x, f)|$ in units of V/V.

\begin{table}[htbp]
\centering
\caption{Complete transfer function magnitudes $|H|$ (V/V) at all 22
distances and 15 frequencies. Values below \num{0.1} are shown with 4
significant figures.}
\label{tab:full-data}
\footnotesize
\setlength{\tabcolsep}{3pt}
\begin{tabular}{@{}r|ccc|ccc|ccc@{}}
\toprule
$x$ (\si{\um}) & 0.1\,Hz & 1\,Hz & 10\,Hz & 50\,Hz & 100\,Hz & 200\,Hz & 500\,Hz & 1000\,Hz & 2000\,Hz \\
\midrule
10   & 0.1607 & 0.1601 & 0.1610 & 0.1605 & 0.1599 & 0.1595 & 0.1526 & 0.1378 & 0.1169 \\
50   & 0.1530 & 0.1530 & 0.1526 & 0.1523 & 0.1518 & 0.1514 & 0.1447 & 0.1304 & 0.1109 \\
100  & 0.1417 & 0.1422 & 0.1421 & 0.1417 & 0.1416 & 0.1411 & 0.1349 & 0.1216 & 0.1040 \\
200  & 0.1255 & 0.1253 & 0.1253 & 0.1253 & 0.1255 & 0.1248 & 0.1191 & 0.1073 & 0.0913 \\
300  & 0.1118 & 0.1117 & 0.1117 & 0.1118 & 0.1119 & 0.1114 & 0.1064 & 0.0960 & 0.0813 \\
500  & 0.0919 & 0.0920 & 0.0919 & 0.0918 & 0.0921 & 0.0917 & 0.0877 & 0.0788 & 0.0671 \\
700  & 0.0779 & 0.0779 & 0.0780 & 0.0782 & 0.0782 & 0.0780 & 0.0744 & 0.0668 & 0.0567 \\
1000 & 0.0636 & 0.0635 & 0.0635 & 0.0637 & 0.0639 & 0.0637 & 0.0607 & 0.0545 & 0.0462 \\
1400 & 0.0509 & 0.0508 & 0.0509 & 0.0512 & 0.0512 & 0.0511 & 0.0487 & 0.0437 & 0.0369 \\
1800 & 0.0424 & 0.0422 & 0.0424 & 0.0426 & 0.0429 & 0.0427 & 0.0405 & 0.0364 & 0.0309 \\
2000 & 0.0390 & 0.0391 & 0.0391 & 0.0392 & 0.0394 & 0.0395 & 0.0374 & 0.0337 & 0.0285 \\
\bottomrule
\end{tabular}

\vspace{0.5em}
\footnotesize
Note: Table shows a selected subset (9 of 15 frequencies, 11 of 22 distances)
for readability. The complete dataset contains $22 \times 15 = 330$ entries,
all valid. Full data available in \texttt{Charge\_Data.csv}.
\end{table}
